\documentclass[11pt,a4paper,openany]{ctexrep}
\usepackage[a4paper,margin=23mm,headheight=15pt]{geometry}
\usepackage{amsmath,amssymb,mathtools,bm}
\usepackage{siunitx}
\usepackage{booktabs,tabularx,array,longtable}
\usepackage{enumitem}
\usepackage{graphicx}
\usepackage{xcolor}
\usepackage{tikz}
\usetikzlibrary{arrows.meta,positioning,calc}
\usepackage[most]{tcolorbox}
\usepackage{hyperref}
\usepackage{fancyhdr}
\usepackage{titlesec}
\usepackage{microtype}
\usepackage{setspace}

\definecolor{KidBlue}{HTML}{22577A}
\definecolor{KidTeal}{HTML}{2A9D8F}
\definecolor{KidGold}{HTML}{E9C46A}
\definecolor{KidRed}{HTML}{C8553D}
\definecolor{KidGray}{HTML}{F3F5F7}
\definecolor{KidDark}{HTML}{263238}

\hypersetup{colorlinks=true,linkcolor=KidBlue,urlcolor=KidBlue,citecolor=KidTeal,pdftitle={KID 器件加工与实验入门},pdfauthor={KID Intro Guide project}}
\setstretch{1.15}
\setlength{\parindent}{2em}
\setlength{\parskip}{0.25em}
\setlist[itemize]{leftmargin=2em,itemsep=0.25em,topsep=0.4em}
\setlist[enumerate]{leftmargin=2em,itemsep=0.25em,topsep=0.4em}

\pagestyle{fancy}
\fancyhf{}
\fancyhead[L]{KID 器件加工与实验入门}
\fancyhead[R]{\leftmark}
\fancyfoot[C]{\thepage}
\renewcommand{\headrulewidth}{0.4pt}

\titleformat{\chapter}[display]{\bfseries\Huge\color{KidDark}}{\filleft\color{KidBlue}\fontsize{44}{44}\selectfont\thechapter}{1ex}{\titlerule\vspace{1ex}\filright}
\titleformat{\section}{\Large\bfseries\color{KidBlue}}{\thesection}{0.7em}{}
\titleformat{\subsection}{\large\bfseries\color{KidTeal}}{\thesubsection}{0.7em}{}

\newtcolorbox{chaptergoal}{colback=KidBlue!5,colframe=KidBlue,boxrule=0.8pt,arc=2mm,title=本章目标,fonttitle=\bfseries}
\newtcolorbox{processbox}{colback=KidTeal!6,colframe=KidTeal,boxrule=0.8pt,arc=2mm,title=工艺 → KID 因果链,fonttitle=\bfseries}
\newtcolorbox{measurebox}{colback=KidGold!16,colframe=KidGold!80!black,boxrule=0.8pt,arc=2mm,title=应该记录什么,fonttitle=\bfseries}
\newtcolorbox{warningbox}{colback=KidRed!5,colframe=KidRed,boxrule=0.8pt,arc=2mm,title=边界与安全,fonttitle=\bfseries}
\newtcolorbox{checkpointbox}{colback=KidGray,colframe=KidDark!55,boxrule=0.7pt,arc=2mm,title=理解检查,fonttitle=\bfseries}

\newcommand{\processchain}[1]{\begin{processbox}#1\end{processbox}}
\newcommand{\kidparam}[1]{\ensuremath{\mathrm{#1}}}
